% GENERATED FILE. Edit canonical lessons, then run npm run generate:books.
% canonical-source-hash: fnv1a64:1b8f10b1969a76da
% canonical-lessons: JA-C484-saihou, JA-C484-amu, JA-C484-keito, JA-C484-hokorobiru, JA-C484-shibaru

\chapter{Needle and Thread}
\label{ch:ja-needle-and-thread}
\hlchaptermodality{484}

\begin{chapteropening}
\textbf{By the end of this chapter:} \emph{I can talk about sewing, knitting and tying.}

Complete the last lesson of chapter 484: I can talk about sewing, knitting and tying.
\end{chapteropening}

\section[\texorpdfstring{\ja{さいほう} (saihō)}{さいほう (saihō)}]{\ja{さいほう} (saihō) — sewing; needlework}

\label{lesson:JA-C484-saihou}



\begin{quote}
Before the new one: say the Japanese for to harden, then the Japanese for to set hard.
\end{quote}

\subsection*{You'll want to know: \ja{さいほう}}
\textbf{\ja{さいほう}} — \emph{saihō} — \textquotedblleft{}sewing; needlework\textquotedblright{}.

\begin{quote}
\textbf{\ja{そぼは} \ja{さいほうが} \ja{じょうずです。}} — \emph{sobo wa saihō ga jōzu desu} — \textquotedblleft{}My grandmother is good at sewing.\textquotedblright{}
\end{quote}

\begin{cousinweb}
\textbf{\ja{さいほう}} (\emph{saihō}) is Sino-Japanese (kango): two Chinese characters read with their on'yomi, readings borrowed from Chinese. They are \emph{sai}, \textquotedblleft{}tailor\textquotedblright{} and \emph{hō}, \textquotedblleft{}sew\textquotedblright{} (Middle Chinese \emph{bjowng}) (KANJIDIC2; Middle Chinese in Baxter's transcription).
\end{cousinweb}

\begin{grammarlens}[title={What you've built}]
One more word for describing what you lived and felt.
\end{grammarlens}

\subsection*{Your turn}
\begin{itemize}
\raggedright
  \item \emph{Say it:} \emph{saihō}
  \item \emph{Say it:} \emph{saihō}, once more
  \item \emph{Say it:} \emph{sobo wa saihō ga jōzu desu} — My grandmother is good at sewing.
  \item \emph{Recall:} say the Japanese for to harden, then the Japanese for to set hard, then say \emph{saihō} again
\end{itemize}

\begin{tcolorbox}[breakable,colback=teal!4,colframe=teal!35!black,title={Before you move on}]
What does \ja{さいほう} mean? (\textquotedblleft{}Sewing; needlework\textquotedblright{}.) Say it once more.
\end{tcolorbox}
\section[\texorpdfstring{\ja{あむ} (amu)}{あむ (amu)}]{\ja{あむ} (amu) — to knit}

\label{lesson:JA-C484-amu}



\begin{quote}
Before the new one: say the Japanese for to set hard, then the Japanese for sewing.
\end{quote}

\subsection*{You'll want to know: \ja{あむ}}
\textbf{\ja{あむ}} — \emph{amu} — \textquotedblleft{}to knit\textquotedblright{}.

\begin{quote}
\textbf{\ja{ははは} \ja{ふゆの} \ja{ぼうしを} \ja{あみます。}} — \emph{haha wa fuyu no bōshi o amimasu} — \textquotedblleft{}My mother knits a winter hat.\textquotedblright{}
\end{quote}

Polite \textbf{\ja{あみます}}, \emph{amimasu}; a godan verb. Past (polite): \textbf{\ja{ははが} \ja{てぶくろを} \ja{あみました。}}, \emph{haha ga tebukuro o amimashita}.

\begin{cousinweb}
\textbf{\ja{あむ}} (\emph{amu}) is a native Japanese word (wago), inherited rather than borrowed from Chinese; its kanji was attached later as a kun'yomi, and its older origin is not settled (JMdict, KANJIDIC2).
\end{cousinweb}

\begin{grammarlens}[title={What you've built}]
One more word for describing what you lived and felt.
\end{grammarlens}

\subsection*{Your turn}
\begin{itemize}
\raggedright
  \item \emph{Say it:} \emph{amu}
  \item \emph{Say it:} \emph{amu}, once more
  \item \emph{Say it:} \emph{haha wa fuyu no bōshi o amimasu} — My mother knits a winter hat.
  \item \emph{Recall:} say the Japanese for to set hard, then the Japanese for sewing, then say \emph{amu} again
\end{itemize}

\begin{tcolorbox}[breakable,colback=teal!4,colframe=teal!35!black,title={Before you move on}]
What does \ja{あむ} mean? (\textquotedblleft{}To knit\textquotedblright{}.) Say it once more.
\end{tcolorbox}
\section[\texorpdfstring{\ja{けいと} (keito)}{けいと (keito)}]{\ja{けいと} (keito) — knitting wool; yarn}

\label{lesson:JA-C484-keito}



\begin{quote}
Before the new one: say the Japanese for sewing, then the Japanese for to knit.
\end{quote}

\subsection*{You'll want to know: \ja{けいと}}
\textbf{\ja{けいと}} — \emph{keito} — \textquotedblleft{}knitting wool; yarn\textquotedblright{}.

\begin{quote}
\textbf{\ja{けいとで} \ja{てぶくろを} \ja{あみます。}} — \emph{keito de tebukuro o amimasu} — \textquotedblleft{}I knit gloves from wool.\textquotedblright{}
\end{quote}

\begin{cousinweb}
\textbf{\ja{けいと}} (\emph{keito}) is a native Japanese compound (wago), inherited rather than borrowed from Chinese: built from \emph{ke}, \textquotedblleft{}hair\textquotedblright{} and \emph{ito}, \textquotedblleft{}thread\textquotedblright{} (JMdict, KANJIDIC2).
\end{cousinweb}

\begin{grammarlens}[title={What you've built}]
One more word for describing what you lived and felt.
\end{grammarlens}

\subsection*{Your turn}
\begin{itemize}
\raggedright
  \item \emph{Say it:} \emph{keito}
  \item \emph{Say it:} \emph{keito}, once more
  \item \emph{Say it:} \emph{keito de tebukuro o amimasu} — I knit gloves from wool.
  \item \emph{Recall:} say the Japanese for sewing, then the Japanese for to knit, then say \emph{keito} again
\end{itemize}

\begin{tcolorbox}[breakable,colback=teal!4,colframe=teal!35!black,title={Before you move on}]
What does \ja{けいと} mean? (\textquotedblleft{}Knitting wool; yarn\textquotedblright{}.) Say it once more.
\end{tcolorbox}
\section[\texorpdfstring{\ja{ほころびる} (hokorobiru)}{ほころびる (hokorobiru)}]{\ja{ほころびる} (hokorobiru) — to come apart at the seams}

\label{lesson:JA-C484-hokorobiru}



\begin{quote}
Before the new one: say the Japanese for to knit, then the Japanese for knitting wool.
\end{quote}

\subsection*{You'll want to know: \ja{ほころびる}}
\textbf{\ja{ほころびる}} — \emph{hokorobiru} — \textquotedblleft{}to come apart at the seams\textquotedblright{}.

\begin{quote}
\textbf{\ja{うわぎの} \ja{そでが} \ja{ほころびました。}} — \emph{uwagi no sode ga hokorobimashita} — \textquotedblleft{}The sleeve of my jacket came apart at the seam.\textquotedblright{}
\end{quote}

Polite \textbf{\ja{ほころびます}}, \emph{hokorobimasu}; an ichidan verb. Past (polite): \textbf{\ja{ふくの} \ja{そでが} \ja{ほころびました。}}, \emph{fuku no sode ga hokorobimashita}.

\begin{cousinweb}
\textbf{\ja{ほころびる}} (\emph{hokorobiru}) is a native Japanese word (wago), inherited rather than borrowed from Chinese, and usually written in kana; its older origin is not settled (JMdict).
\end{cousinweb}

\begin{grammarlens}[title={What you've built}]
One more word for describing what you lived and felt.
\end{grammarlens}

\subsection*{Your turn}
\begin{itemize}
\raggedright
  \item \emph{Say it:} \emph{hokorobiru}
  \item \emph{Say it:} \emph{hokorobiru}, once more
  \item \emph{Say it:} \emph{uwagi no sode ga hokorobimashita} — The sleeve of my jacket came apart at the seam.
  \item \emph{Recall:} say the Japanese for to knit, then the Japanese for knitting wool, then say \emph{hokorobiru} again
\end{itemize}

\begin{tcolorbox}[breakable,colback=teal!4,colframe=teal!35!black,title={Before you move on}]
What does \ja{ほころびる} mean? (\textquotedblleft{}To come apart at the seams\textquotedblright{}.) Say it once more.
\end{tcolorbox}
\section[\texorpdfstring{\ja{しばる} (shibaru)}{しばる (shibaru)}]{\ja{しばる} (shibaru) — to tie up; to bind}

\label{lesson:JA-C484-shibaru}



\begin{quote}
Before the new one: say the Japanese for knitting wool, then the Japanese for to come apart at the seams.
\end{quote}

\subsection*{You'll want to know: \ja{しばる}}
\textbf{\ja{しばる}} — \emph{shibaru} — \textquotedblleft{}to tie up; to bind\textquotedblright{}.

\begin{quote}
\textbf{\ja{ふるい} \ja{しんぶんを} \ja{ひもで} \ja{しばります。}} — \emph{furui shinbun o himo de shibarimasu} — \textquotedblleft{}I tie up the old newspapers with string.\textquotedblright{}
\end{quote}

Polite \textbf{\ja{しばります}}, \emph{shibarimasu}; a godan verb. Past (polite): \textbf{\ja{ちちが} \ja{にもつを} \ja{しばりました。}}, \emph{chichi ga nimotsu o shibarimashita}.

\begin{cousinweb}
\textbf{\ja{しばる}} (\emph{shibaru}) is a native Japanese word (wago), inherited rather than borrowed from Chinese; its kanji was attached later as a kun'yomi, and its older origin is not settled (JMdict, KANJIDIC2).
\end{cousinweb}

\begin{grammarlens}[title={What you've built}]
One more word for describing what you lived and felt.
\end{grammarlens}

\subsection*{Your turn}
\begin{itemize}
\raggedright
  \item \emph{Say it:} \emph{shibaru}
  \item \emph{Say it:} \emph{shibaru}, once more
  \item \emph{Say it:} \emph{furui shinbun o himo de shibarimasu} — I tie up the old newspapers with string.
  \item \emph{Recall:} say the Japanese for knitting wool, then the Japanese for to come apart at the seams, then say \emph{shibaru} again
\end{itemize}

\begin{tcolorbox}[breakable,colback=teal!4,colframe=teal!35!black,title={Before you move on}]
What does \ja{しばる} mean? (\textquotedblleft{}To tie up; to bind\textquotedblright{}.) Say it once more.
\end{tcolorbox}
